% GENERATED FILE. Edit canonical lessons, then run npm run generate:books.
% canonical-source-hash: fnv1a64:d25eecf4715b68c5
% canonical-lessons: KA-C186-sandesa, KA-C186-imel, KA-C186-pensil, KA-C186-kadata, KA-C186-dakhale

\chapter{Message, Email, Pencil, File, Document}
\label{ch:ka-message-email-pencil-file-document}
\hlchaptermodality{186}

\begin{chapteropening}
\textbf{By the end of this chapter:} \emph{I can say a message, an email, a pencil, a file and a document.}

Complete the last lesson of chapter 186: I can say a message, an email, a pencil, a file and a document.
\end{chapteropening}

\section[\texorpdfstring{\kn{ಸಂದೇಶ} (sandēśa)}{ಸಂದೇಶ (sandēśa)}]{\kn{ಸಂದೇಶ} (sandēśa) — a message}

\label{lesson:KA-C186-sandesa}



\begin{quote}
Before the new one: say the Kannada for a kilogram, then the Kannada for a purse.
\end{quote}

\subsection*{You'll want to know: \kn{ಸಂದೇಶ}}
\textbf{\kn{ಸಂದೇಶ}} — \emph{sandēśa} — \textquotedblleft{}a message\textquotedblright{}.

\begin{grammarlens}[title={What you've built}]
One more word for this chapter.
\end{grammarlens}

\subsection*{Your turn}
\begin{itemize}
\raggedright
  \item \emph{Say it:} \emph{sandēśa}
  \item \emph{Say it:} \emph{sandēśa}, once more
  \item \emph{Say it:} say \emph{sandēśa}
  \item \emph{Recall:} say the Kannada for a kilogram, then the Kannada for a purse, then say \emph{sandēśa} again
\end{itemize}

\begin{tcolorbox}[breakable,colback=teal!4,colframe=teal!35!black,title={Before you move on}]
What does \kn{ಸಂದೇಶ} mean? (\textquotedblleft{}A message\textquotedblright{}.) Say it once more.
\end{tcolorbox}
\section[\texorpdfstring{\kn{ಇಮೇಲ್} (imēl)}{ಇಮೇಲ್ (imēl)}]{\kn{ಇಮೇಲ್} (imēl) — an email}

\label{lesson:KA-C186-imel}



\begin{quote}
Before the new one: say the Kannada for a purse, then the Kannada for a message.
\end{quote}

\subsection*{You'll want to know: \kn{ಇಮೇಲ್}}
\textbf{\kn{ಇಮೇಲ್}} — \emph{imēl} — \textquotedblleft{}an email\textquotedblright{}.

\begin{grammarlens}[title={What you've built}]
One more word for this chapter.
\end{grammarlens}

\subsection*{Your turn}
\begin{itemize}
\raggedright
  \item \emph{Say it:} \emph{imēl}
  \item \emph{Say it:} \emph{imēl}, once more
  \item \emph{Say it:} say \emph{imēl}
  \item \emph{Recall:} say the Kannada for a purse, then the Kannada for a message, then say \emph{imēl} again
\end{itemize}

\begin{tcolorbox}[breakable,colback=teal!4,colframe=teal!35!black,title={Before you move on}]
What does \kn{ಇಮೇಲ್} mean? (\textquotedblleft{}An email\textquotedblright{}.) Say it once more.
\end{tcolorbox}
\section[\texorpdfstring{\kn{ಪೆನ್ಸಿಲ್} (pensil)}{ಪೆನ್ಸಿಲ್ (pensil)}]{\kn{ಪೆನ್ಸಿಲ್} (pensil) — a pencil}

\label{lesson:KA-C186-pensil}



\begin{quote}
Before the new one: say the Kannada for a message, then the Kannada for an email.
\end{quote}

\subsection*{You'll want to know: \kn{ಪೆನ್ಸಿಲ್}}
\textbf{\kn{ಪೆನ್ಸಿಲ್}} — \emph{pensil} — \textquotedblleft{}a pencil\textquotedblright{}.

\begin{grammarlens}[title={What you've built}]
One more word for this chapter.
\end{grammarlens}

\subsection*{Your turn}
\begin{itemize}
\raggedright
  \item \emph{Say it:} \emph{pensil}
  \item \emph{Say it:} \emph{pensil}, once more
  \item \emph{Say it:} say \emph{pensil}
  \item \emph{Recall:} say the Kannada for a message, then the Kannada for an email, then say \emph{pensil} again
\end{itemize}

\begin{tcolorbox}[breakable,colback=teal!4,colframe=teal!35!black,title={Before you move on}]
What does \kn{ಪೆನ್ಸಿಲ್} mean? (\textquotedblleft{}A pencil\textquotedblright{}.) Say it once more.
\end{tcolorbox}
\section[\texorpdfstring{\kn{ಕಡತ} (kaḍata)}{ಕಡತ (kaḍata)}]{\kn{ಕಡತ} (kaḍata) — a file (of papers)}

\label{lesson:KA-C186-kadata}



\begin{quote}
Before the new one: say the Kannada for an email, then the Kannada for a pencil.
\end{quote}

\subsection*{You'll want to know: \kn{ಕಡತ}}
\textbf{\kn{ಕಡತ}} — \emph{kaḍata} — \textquotedblleft{}a file (of papers)\textquotedblright{}.

\begin{grammarlens}[title={What you've built}]
One more word for this chapter.
\end{grammarlens}

\subsection*{Your turn}
\begin{itemize}
\raggedright
  \item \emph{Say it:} \emph{kaḍata}
  \item \emph{Say it:} \emph{kaḍata}, once more
  \item \emph{Say it:} say \emph{kaḍata}
  \item \emph{Recall:} say the Kannada for an email, then the Kannada for a pencil, then say \emph{kaḍata} again
\end{itemize}

\begin{tcolorbox}[breakable,colback=teal!4,colframe=teal!35!black,title={Before you move on}]
What does \kn{ಕಡತ} mean? (\textquotedblleft{}A file (of papers)\textquotedblright{}.) Say it once more.
\end{tcolorbox}
\section[\texorpdfstring{\kn{ದಾಖಲೆ} (dākhale)}{ದಾಖಲೆ (dākhale)}]{\kn{ದಾಖಲೆ} (dākhale) — a document, a record}

\label{lesson:KA-C186-dakhale}



\begin{quote}
Before the new one: say the Kannada for a pencil, then the Kannada for a file.
\end{quote}

\subsection*{You'll want to know: \kn{ದಾಖಲೆ}}
\textbf{\kn{ದಾಖಲೆ}} — \emph{dākhale} — \textquotedblleft{}a document, a record\textquotedblright{}.

\begin{grammarlens}[title={What you've built}]
One more word for this chapter.
\end{grammarlens}

\subsection*{Your turn}
\begin{itemize}
\raggedright
  \item \emph{Say it:} \emph{dākhale}
  \item \emph{Say it:} \emph{dākhale}, once more
  \item \emph{Say it:} say \emph{dākhale}
  \item \emph{Recall:} say the Kannada for a pencil, then the Kannada for a file, then say \emph{dākhale} again
\end{itemize}

\begin{tcolorbox}[breakable,colback=teal!4,colframe=teal!35!black,title={Before you move on}]
What does \kn{ದಾಖಲೆ} mean? (\textquotedblleft{}A document, a record\textquotedblright{}.) Say it once more.
\end{tcolorbox}
